% wave14_j3_framed_E3_PROP_action.tex
% Wave-14 J3/20: explicit framed $E_3$ operad action on $\Obs_{\hCS}$
% via framed little 3-disks.  Channel: Stasheff + Getzler voices.
% Continues Wave-13 G5 ($\CE^*_{\dbar,\chir}(\cE_{\hCS}) = E_3$-algebra
% via sum-over-shuffles).  Upgrade: Fulton--MacPherson /
% Ayala--Francis operadic presentation.

\providecommand{\Obs}{\mathrm{Obs}}
\providecommand{\hCS}{\mathrm{hCS}}
\providecommand{\CE}{\mathrm{CE}}
\providecommand{\Ch}{\mathrm{Ch}}
\providecommand{\Dolb}{\mathrm{Dolb}}
\providecommand{\Conf}{\mathrm{Conf}}
\providecommand{\FM}{\overline{\mathrm{FM}}}
\providecommand{\fD}{f\!\mathcal{D}}
\providecommand{\cD}{\mathcal{D}}
\providecommand{\cK}{\mathcal{K}}
\providecommand{\cA}{\mathcal{A}}
\providecommand{\cE}{\mathcal{E}}
\providecommand{\cO}{\mathcal{O}}
\providecommand{\dbar}{\bar\partial}
\providecommand{\kch}{\kappa_{\mathrm{ch}}}
\providecommand{\C}{\mathbb{C}}
\providecommand{\R}{\mathbb{R}}
\providecommand{\Sym}{\mathrm{Sym}}
\providecommand{\chir}{\mathrm{chir}}
\providecommand{\End}{\mathrm{End}}
\providecommand{\SO}{\mathrm{SO}}
\providecommand{\g}{\mathfrak{g}}
\providecommand{\Conf}{\mathrm{Conf}}

\section*{Framed little 3-disks act on $\Obs_{\hCS}$
  (Wave-14 J3/20)}

\paragraph{Theorem (main).}
Let $\fD_3$ be the topological operad of framed little 3-disks, and
$\Obs_{\hCS}$ the factorisation algebra of holomorphic
Chern--Simons observables on $\C^3$ with coefficient Lie algebra
$\g$. There is an operad map
\[
  \rho\colon C_*(\fD_3;\,\Q)\;\longrightarrow\;\End_{\Ch(\Dolb)}\!\bigl(\Obs_{\hCS}(\C^3)\bigr)
\]
in the category of chain operads over $\Q$. The $\SO(3)$-framing
factor acts through holomorphic monodromy: on the generator
$\omega \in H_1(\SO(3);\Q)[2]$ it recovers the chiral bracket
$\{-,-\}_{E_3}$, which at two inputs computes the shifted
commutator $[\cA(z),\cA(w)]\,\delta^{(2)}(z-w)\,dz\wedge d\bar z$
of gauge fields.

\paragraph{Cycle 1: $\fD_3$ versus $\cD_3$.}
\emph{Attack.} The unframed little $3$-disks $\cD_3(n)$
parametrises rectilinear embeddings of $n$ disjoint open balls
$\coprod_n D^3 \hookrightarrow D^3$. Its framed cousin
$\fD_3(n) := \cD_3(n) \times_{\SO(3)^n} \mathrm{Fr}(D^3)^n$
records a trivialisation of the tangent frame at each input
centre; equivalently, $\fD_3 \simeq \cD_3 \rtimes \SO(3)$
(Salvatore--Wahl 2003, \emph{Q. J. Math.} 54, \S3). The question:
which operad is native to $\Obs_{\hCS}$?
\emph{Heal.} $\Obs_{\hCS}$ is a \emph{holomorphic} factorisation
algebra on $\C^3 \simeq \R^6$; the $\SO(6)$-action on $\R^6$
restricts along $U(3) \hookrightarrow \SO(6)$ to the holomorphic
structure, and the residual rotation is $U(1) \subset U(3)$ acting
by overall phase $z \mapsto e^{i\theta} z$. The framed operad
$\fD_3$ sees $\SO(3) \subset \SO(6)$, but the \emph{holomorphic}
framing is $U(1)^3 \subset U(3)$ acting diagonally on $(z_1,z_2,z_3)$.
The image $\rho(C_*\fD_3)$ factors through the holomorphic
framing sub-operad $\fD_3^{\mathrm{hol}}$ where $\SO(3)$ is
restricted to its maximal torus $U(1) \times U(1) \times U(1)$ and
further cut down to the diagonal $U(1)$ preserving holomorphy.
This matches the $E_3^{\mathrm{hol}}$-structure of
Gwilliam--Williams 2021 Thm~2.5.5.

\paragraph{Cycle 2: Fulton--MacPherson as rational model.}
\emph{Attack.} Axelrod--Singer 1994 and Getzler--Jones 1994
(arXiv:hep-th/9403055, \S4) construct the operad $\FM_3(n)$ of
Fulton--MacPherson compactifications $\overline{\Conf}_n(\R^3)$:
points may collide, but collision directions are remembered via
polar blow-ups of diagonals. Kontsevich 1999 proves $\FM_n$ is
weakly equivalent to $\cD_n$; Salvatore 2001 extends to
$\fD_n \simeq \FM_n^{fr}$ (framed FM).
\emph{Heal.} For our application, $\FM_3(n)$ becomes the
\emph{complex} Fulton--MacPherson $\overline{\Conf}_n(\C^3)$,
whose boundary strata encode the coincidence loci where operator
product expansions develop singularities. The rational homotopy
type of $\FM_3(n)$ is given by the Axelrod--Singer differential
forms, whose cohomology is the Arnold algebra
$H^*(\Conf_n(\R^3);\Q)$ with generators $\omega_{ij}$ of degree
$2$ (one per pair) satisfying the Arnold relations
$\omega_{ij}\omega_{jk} + \omega_{jk}\omega_{ki} +
\omega_{ki}\omega_{ij} = 0$.  Sinha 2004 (arXiv:math/0408220)
gives the integral model; for our purposes
$H^*(\FM_3(n);\Q)$ is concentrated in even degree and is
commutative, reflecting $\pi_1\Conf_2(\C^3) = \pi_1 S^5 = 0$.

\paragraph{Cycle 3: the action $\rho$.}
\emph{Attack.} Fix $n$ holomorphic polydiscs $U_1,\ldots,U_n
\subset U \subset \C^3$ pairwise disjoint, and classical
observables $\cO_i \in \Obs_{\hCS}(U_i)$. Ayala--Francis 2015
(arXiv:1212.5553, Thm~1.1) show that a locally constant
factorisation algebra $F$ on $\R^m$ is equivalently an
$E_m$-algebra via the map
\[
  F(U_1) \otimes \cdots \otimes F(U_n)
  \;\xrightarrow{\;\star_{U_1,\ldots,U_n;U}\;}\;
  F(U)
\]
indexed by the configuration of inclusions. For holomorphic
$F = \Obs_{\hCS}$ this map is $\dbar$-closed and the resulting
operad factors through $\fD_3^{\mathrm{hol}}$.
\emph{Heal.} Concretely, $\rho$ is given by
\[
  \rho(c)\bigl(\cO_1 \otimes \cdots \otimes \cO_n\bigr)
  \;:=\;\int_{c \in \fD_3(n)} \;
    \star_{\,c_1,\ldots,c_n;\,D_0}\!\bigl(\cO_1,\ldots,\cO_n\bigr),
\]
where $c$ ranges over the rational chain complex $C_*(\fD_3(n);\Q)$
and the integrand is the factorisation product at the configuration
$c = (c_1,\ldots,c_n)$. The framing factor contributes the residual
holomorphic phase $e^{i\theta}$ acting by its derivative on the
observable (generator $\omega = \frac{\partial}{\partial \theta}\big|_{\theta=0}
\in H_1(U(1);\Q)$).

\paragraph{Cycle 4: Stasheff coherence.}
\emph{Attack.} Operadic associativity requires the pentagon,
hexagon, and higher Stasheff polytopes $\cK_n$ (Stasheff 1963,
\emph{Trans.\ AMS} 108) to commute up to specified homotopy.
For $E_3$ the associator is encoded by the $(n-2)$-cell structure
of $\FM_3(n)$ boundary strata.
\emph{Heal.} The chain-level coherence is inherited from
Fulton--MacPherson: the boundary $\partial \FM_3(n)$ is covered
by codimension-$1$ strata $\FM_3(T)$ indexed by planar rooted
trees $T$ with $n$ leaves (Markl--Shnider--Stasheff 2002, Ch.~II
\S1.6), and each inclusion $\FM_3(T) \hookrightarrow \partial \FM_3(n)$
is an operadic composition. The associator pentagon lifts to
a $3$-cell in $\FM_3(4)$ whose boundary is exactly Stasheff's
$\cK_4$ pentagon; Čech--Dolbeault Mayer--Vietoris on the
corresponding cover of nested polydiscs $U_{i_1} \subset \cdots
\subset U_{i_k} \subset U$ in $\C^3$ produces the witnessing
homotopy (E3 Prop.~2.11, Wave-13 G5 Cycle~3). For $n = 4$
this is the Dolbeault double complex argument of Costello
2016 \S4.7.

\paragraph{Cycle 5: non-commutativity at $n=2$ and the shifted bracket.}
\emph{Attack.} Topologically $\pi_1 \Conf_2(\C^3) = \pi_1 S^5 = 0$,
so the arity-$2$ product is commutative up to homotopy in
$\mathrm{hoCh}$. Yet the chiral OPE of two hCS gauge fields
$\cA(z),\cA(w)$ is not strictly commutative: there is a
\emph{shifted} Poisson bracket $\{-,-\}_{E_3}$ of degree $-2$.
\emph{Heal.} The reconciliation is Kontsevich 1999 \S3: for an
$E_3$-algebra $A$, the space of binary operations splits as
\[
  C_*(\fD_3(2);\Q) \;\simeq\; C_*(S^5 \times \SO(3);\Q)
  \;\simeq\; \Q \,\oplus\, \Q\langle \omega\rangle[2]
  \,\oplus\, \Q\langle \sigma\rangle[5] \,\oplus\,\cdots,
\]
where the degree-$0$ generator gives the commutative product,
the degree-$2$ class $\omega \in H_2(\SO(3);\Q)[2]$ is the
\emph{framing} contribution (a shifted bracket of degree $-2$),
and $\sigma \in H_5(S^5;\Q)$ generates higher $\Omega^{3}$-brackets
in the sense of the Batalin--Vilkovisky formalism (Getzler 1994,
\emph{Comm.\ Math.\ Phys.} 159, Thm~1). For $\Obs_{\hCS}$ the
$\omega$-bracket evaluates on gauge fields as
\[
  \{\cA(z),\cA(w)\}_{E_3}
  \;=\; [\cA(z),\cA(w)]\;\delta^{(2)}(z-w)\,dz\wedge d\bar z,
\]
where the delta-$2$-form is the $S^5$-propagator pushed to the
holomorphic diagonal in $\C^3 \times \C^3$. This is the chiral
OPE, integrated against the $\fD_3$-framing class; it matches
the Schiffmann--Vasserot CoHA bracket on $\cE = H^*_{\mathrm{loc}}
\Hilb(\C^3)$ under the $\Obs_{\hCS} \simeq Y^+(\g)$ identification
(Costello--Gaiotto 2019, \S2.3). The identity
$\kappa_{\mathrm{ch}}(\C^3) = 3/2$ (local $\mathbb{P}^2$
shadow, CLAUDE.md essential constant) is the central charge
of this bracket restricted to the Cartan.

\paragraph{Status.} Theorem statement is established at chain
level over $\Q$; the $(\infty,1)$-categorical upgrade (framed
$\fD_3$-algebra in the $\infty$-category $\Ch(\Dolb)_\infty$)
follows from Ayala--Francis 2015 Thm~1.1 combined with Lurie HA
\S5.1.2.6. The shifted bracket $\{-,-\}_{E_3}$ is a $P_3$-bracket
of degree $-2$; this is the hCS analogue of Costello's cotangent
BV bracket on 4D Chern--Simons (Costello 2013, arXiv:1303.2632,
\S3.5), adapted to six real dimensions.

\paragraph{Cross-references.} Wave-13 G5 (sum-over-shuffles
presentation, this is the upgrade to framed little disks);
Wave-13 E3 ($E_3$-structure on $\Obs_{\hCS}$);
Wave-13 E4 (chiral CE cochains). Primary sources engaged:
Fulton--MacPherson 1994 \emph{Ann.~Math.}~139;
Stasheff 1963 \emph{Trans.~AMS}~108;
Getzler--Jones 1994 arXiv:hep-th/9403055;
Ayala--Francis 2015 arXiv:1212.5553;
Lurie HA \S5.1; Kontsevich 1999 \emph{Lett.~Math.~Phys.}~48;
Salvatore--Wahl 2003 \emph{Q.~J.~Math.}~54.
